\chapter{透镜、像与眼睛}\label{ch:g11:lenses-and-eye}

清晰的像——无论出现在手机感光元件上、电影银幕上，还是
你自己眼睛的后部——都源于同一个技巧：弯曲的透明介质借助折射
（见\cref{ch:g10:refraction}）使光线偏折，使所有从同一点射出的光线
再次会聚于另一点。本章驾驭这一技巧：
一个透镜、三条光线、一个公式——然后将它用于
相机、眼镜，以及正在阅读这句话的活的晶状体。

\section{凸透镜与凹透镜}

\begin{definition}[薄透镜]\label{def:g11:lenses-and-eye:lens}
由两个表面围成的透明介质称为\emph{透镜}\index{透镜}，其中至少
有一个表面是弯曲的。中央比边缘厚的透镜
称为\emph{凸透镜}\index{凸透镜}：平行光线被
弯折后相互靠拢；边缘比中央厚的透镜称为\emph{凹透镜}\index{凹透镜}：
平行光线向外散开。\emph{薄透镜}\index{薄透镜}的厚度可忽略，
画成线段，两端带向外（凸透镜）或
向内（凹透镜）的箭头。其中心 $O$ 为
\emph{光心}\index{光心}；过 $O$ 且垂直于透镜的直线，即
\emph{主光轴}\index{主光轴}。
\end{definition}

\begin{definition}[焦点与焦距]\label{def:g11:lenses-and-eye:focal}
平行于凸透镜主光轴入射的光线，经折射后都穿过
透镜后方主光轴上的同一点，称为\emph{像方焦点}\index{焦点} $F'$；对称地，
从\emph{物方焦点} $F$（位于透镜前方，且 $OF = OF'$）出发的光线
经折射后平行于主光轴射出。
\emph{焦距}\index{焦距} $f' = \overline{OF'}$
对凸透镜为正；对凹透镜，平行光线经折射后
如同从透镜\emph{前方}的一点 $F'$ 发出，$f' < 0$。
\end{definition}

\begin{definition}[光焦度]\label{def:g11:lenses-and-eye:vergence}
焦距为 $f'$ 的透镜，其\emph{光焦度}\index{光焦度}
（以米计）为 $C = 1/f'$，
单位为\emph{屈光度}\index{屈光度}（符号 D）：
$1~\text{D} = \qty{1}{m^{-1}}$；凹透镜的 $C < 0$。
验光师以屈光度开具处方：相互紧贴的薄透镜的光焦度
直接相加（在本阶段不加证明地承认；推导见本科一年级分册）。
\end{definition}

\begin{example}[读懂验光处方]\label{ex:g11:lenses-and-eye:vergence}
焦距为 $f' = \qty{0.50}{m}$ 的透镜，其 $C = 1/0.50 = +2.0$~D，
为凸透镜。一张 $-4.0$~D 的处方表示 $f' = \qty{-0.25}{m}$，
为凹透镜，焦距为 \qty{25}{cm}。
\end{example}

\section{作图求像}

\begin{method}[作图的三条光线]\label{met:g11:lenses-and-eye:rays}
设 $AB$ 为在 $A$ 处垂直于主光轴的物体。在从端点
$B$ 发出的所有光线中，有三条的路径是已知的：
\begin{enumerate}
  \item 过光心 $O$ 的光线沿直线传播；
  \item 平行于主光轴的光线经 $F'$ 射出；
  \item 过 $F$ 的光线经折射后平行于主光轴射出。
\end{enumerate}
任取其中两条即可：它们的交点就是像点 $B'$，并且
像 $A'B'$ 垂直于主光轴，位于 $A'$ 处。若出射光线
发散，则其\emph{反向}延长线（虚线）相交于
一个虚像点。
\end{method}

\begin{omfigure}
\begin{tikzpicture}[scale=0.55]
  \draw[-Stealth] (-6.4,0) -- (6.4,0);
  \draw[Stealth-Stealth, omDef, very thick] (0,-2.4) -- (0,2.4);
  \fill (0,0) circle (1.5pt) node[below right, font=\small] {$O$};
  \fill (-2,0) circle (1.5pt) node[below left, font=\small] {$F$};
  \fill (2,0) circle (1.5pt) node[below right, font=\small] {$F'$};
  \draw[omMeth] (-4,1.2) -- (0,1.2) -- (4,-1.2);
  \draw[omMeth] (-4,1.2) -- (4,-1.2);
  \draw[omMeth] (-4,1.2) -- (0,-1.2) -- (4,-1.2);
  \draw[-Stealth, very thick] (-4,0) -- (-4,1.2);
  \node[above, font=\small] at (-4,1.2) {$B$};
  \node[below, font=\small] at (-4,0) {$A$};
  \draw[-Stealth, omThm, very thick] (4,0) -- (4,-1.2);
  \node[below, font=\small] at (4,-1.2) {$B'$};
  \node[above, font=\small] at (4,0) {$A'$};
\end{tikzpicture}

{\small 三条作图光线：物 $AB$ 位于 $F$ 外侧，像
$A'B'$ 为实像且倒立（此时 $\overline{OA} = -2f'$，等大，
$\gamma = -1$）。}
\end{omfigure}

\begin{definition}[实像与虚像]\label{def:g11:lenses-and-eye:images}
出射光线确实经过的像称为\emph{实像}\index{实像}，
此时可在该处放置屏幕接收它；
若出射光线只是其反向延长线相交，则称为\emph{虚像}\index{虚像}，
此时屏幕无法接收它，但透过透镜注视的眼睛仍能
清楚地看到它。凸透镜对任一位于 $F$ 以外的物体成倒立的实像，
若将物体移至焦距以内，
则像变为正立、放大的虚像。
\end{definition}

\begin{omfigure}
\begin{tikzpicture}[scale=0.5]
  \draw[-Stealth] (-7.2,0) -- (4.6,0);
  \draw[Stealth-Stealth, omDef, very thick] (0,-2.2) -- (0,3.6);
  \fill (0,0) circle (1.5pt) node[below right, font=\small] {$O$};
  \fill (-2,0) circle (1.5pt) node[below, font=\small] {$F$};
  \fill (2,0) circle (1.5pt) node[below, font=\small] {$F'$};
  \draw[omMeth] (-1.5,0.8) -- (3,-1.6);
  \draw[omMeth] (-1.5,0.8) -- (0,0.8) -- (4,-0.8);
  \draw[omMeth, dashed] (0,0.8) -- (-6,3.2);
  \draw[omMeth, dashed] (-1.5,0.8) -- (-6,3.2);
  \draw[-Stealth, very thick] (-1.5,0) -- (-1.5,0.8);
  \node[below, font=\small] at (-1.5,0) {$A$};
  \node[above right, font=\small] at (-1.4,0.75) {$B$};
  \draw[-Stealth, omThm, very thick, dashed] (-6,0) -- (-6,3.2);
  \node[below, font=\small] at (-6,0) {$A'$};
  \node[above, font=\small] at (-6,3.2) {$B'$};
\end{tikzpicture}

{\small 位于焦距之内的物体（放大镜的设置）：
反向延长线构成一个虚的、正立的、放大的像 $A'B'$。}
\end{omfigure}

\section{一个公式：共轭关系}

作图给出像\emph{在何处}；公式则可以计算它。
沿光的传播方向取坐标轴，并使用\emph{代数量}：
$\overline{OA}$ 为 $A$ 相对于 $O$ 的坐标（对于
前方的实物取负值）；高度 $\overline{AB}$ 向上取正。

\begin{theorem}[薄透镜共轭关系]\label{thm:g11:lenses-and-eye:conjugate}
对于焦距为 $f'$ 的薄透镜，位于主光轴上的物点 $A$
经透镜成像于点 $A'$，且
\[
\frac{1}{\overline{OA'}} - \frac{1}{\overline{OA}} = \frac{1}{f'} = C .
\]
\end{theorem}
\admitted

\begin{remark}
该关系可由斯涅耳定律应用于近轴光线推出；
推导过程以及薄透镜模型的适用范围——
将在本科一年级分册中严格说明。
\end{remark}

\begin{definition}[放大率]\label{def:g11:lenses-and-eye:magnification}
像的\emph{放大率}\index{放大率}为
\[
\gamma = \frac{\overline{A'B'}}{\overline{AB}}
       = \frac{\overline{OA'}}{\overline{OA}} .
\]
$\abs{\gamma}$ 比较大小，其符号表示方向（$\gamma < 0$：
倒立）；$\overline{OA'}$ 的符号反映像的性质（正：
实像，位于透镜后方；负：虚像，位于前方）。
\end{definition}

\begin{example}[公式的应用]\label{ex:g11:lenses-and-eye:conjugate}
物体位于凸透镜前 \qty{30}{cm} 处，焦距为
$f' = \qty{10}{cm}$：$\overline{OA} = \qty{-30}{cm}$，于是
\[
\frac{1}{\overline{OA'}} = \frac{1}{10} - \frac{1}{30} = \frac{1}{15},
\qquad \overline{OA'} = +\qty{15}{cm},
\qquad \gamma = \frac{15}{-30} = -0.50 :
\]
实像位于透镜后 \qty{15}{cm} 处，倒立，大小为物的一半——
正是三条光线所画出的图样。
\end{example}

\section{眼睛}

\begin{definition}[作为光学系统的眼睛]\label{def:g11:lenses-and-eye:eye}
从光学角度看，眼睛是一个正对屏幕的凸透镜：
\emph{角膜}\index{角膜}与\emph{晶状体}\index{晶状体}
合起来相当于一个焦距可调的凸透镜；
\emph{视网膜}\index{视网膜}位于其后约 \qty{17}{mm} 处，是固定的
屏幕。对焦靠改变透镜的曲度来实现，而非像
照相机那样移动透镜：\emph{调节}\index{调节}是由睫状肌引起的
$f'$ 的改变。在完全调节时能清晰看到的最近点
称为\emph{近点}\index{近点}，
对于\emph{标准眼}\index{标准眼}约为 \qty{25}{cm}；
静止时能清晰看到的最远点为\emph{远点}\index{远点}（对于
标准眼为无穷远）。
\end{definition}

\begin{omfigure}
\begin{tikzpicture}[scale=0.85]
  \draw[thick] (0,0) circle (1.1);
  \draw[omDef, very thick] (-0.95,-0.5) -- (-0.95,0.5);
  \draw (-2.6,0.35) -- (-0.95,0.35) -- (1.1,0);
  \draw (-2.6,0) -- (1.1,0);
  \draw (-2.6,-0.35) -- (-0.95,-0.35) -- (1.1,0);
  \fill (1.1,0) circle (1.5pt) node[right, font=\small] {视网膜};
  \node[font=\small] at (-0.7,-1.5) {静止：远处物体};
\end{tikzpicture}\qquad
\begin{tikzpicture}[scale=0.85]
  \draw[thick] (0,0) circle (1.1);
  \draw[omDef, line width=2.6pt] (-0.95,-0.5) -- (-0.95,0.5);
  \fill (-2.6,0) circle (1.5pt) node[above, font=\small] {$B$};
  \draw (-2.6,0) -- (-0.95,0.4) -- (1.1,0);
  \draw (-2.6,0) -- (1.1,0);
  \draw (-2.6,0) -- (-0.95,-0.4) -- (1.1,0);
  \node[font=\small] at (-0.7,-1.5) {调节后：近处物体};
\end{tikzpicture}

{\small 调节：为使近处物体仍清晰地成像于同一视网膜上，
晶状体鼓起——其光焦度增大。}
\end{omfigure}

\begin{example}[你眼睛的光焦度]\label{ex:g11:lenses-and-eye:eyerange}
视网膜位于 $\overline{OA'} = +\qty{0.017}{m}$ 处。远处物体
（$1/\overline{OA} \approx 0$）：$C = 1/0.017 \approx 59$~D。近点
（$\overline{OA} = \qty{-0.25}{m}$）：$C = 1/0.017 + 1/0.25 \approx
63$~D。阅读约需 $4$~D：眼睛相当于一个 $59$~D 的透镜，并带有
内置的 $+4$~D 微调。
\end{example}

\section{缺陷、矫正与放大镜}

\begin{definition}[近视与远视]\label{def:g11:lenses-and-eye:defects}
\emph{近视}\index{近视}（短视）的眼睛会聚过强，
相对于其长度：平行光线聚焦于视网膜\emph{前方}，远处
物体模糊，远点位于有限距离处；可用
凹透镜式眼镜片加以矫正。\emph{远视}\index{远视}
（远视眼）的会聚能力不足：近处物体将聚焦于
视网膜后方，近点向远处移动；需用凸透镜加以矫正。
衰老使晶状体变硬
（\emph{老花眼}\index{老花眼}），调节能力减弱——
矫正方法相同：老花镜。
\end{definition}

\begin{example}[为近视者配镜]\label{ex:g11:lenses-and-eye:correction}
近视眼的远点在 \qty{50}{cm} 处，矫正透镜
须使远处物体成像于\emph{远点}处——物体位于
无穷远，像位于 $\overline{OA'} = \qty{-0.50}{m}$（虚像，位于
镜前）。于是 $1/f' = 1/\overline{OA'}$，$f' = \qty{-0.50}{m}$，
$C = -2.0$~D；松弛状态下的眼睛注视该虚像，
即可清晰地看到它。
\end{example}

\begin{omfigure}
\begin{tikzpicture}[scale=0.85]
  \draw[thick] (0,0) circle (1.1);
  \draw[omDef, very thick] (-0.95,-0.5) -- (-0.95,0.5);
  \draw (-2.7,0.35) -- (-0.95,0.35) -- (0.55,0) -- (1.1,-0.13);
  \draw (-2.7,0) -- (1.1,0);
  \draw (-2.7,-0.35) -- (-0.95,-0.35) -- (0.55,0) -- (1.1,0.13);
  \fill[omThm] (0.55,0) circle (1.5pt);
  \node[font=\small] at (-0.7,-1.5) {近视：焦点位于视网膜之前};
\end{tikzpicture}\qquad
\begin{tikzpicture}[scale=0.85]
  \draw[thick] (0,0) circle (1.1);
  \draw[omDef, very thick] (-0.95,-0.5) -- (-0.95,0.5);
  \draw[omMeth, very thick, {Stealth[reversed]}-{Stealth[reversed]}]
    (-1.9,-0.55) -- (-1.9,0.55);
  \draw (-2.8,0.3) -- (-1.9,0.3) -- (-0.95,0.42) -- (1.1,0);
  \draw (-2.8,0) -- (1.1,0);
  \draw (-2.8,-0.3) -- (-1.9,-0.3) -- (-0.95,-0.42) -- (1.1,0);
  \node[font=\small] at (-0.7,-1.5) {矫正后：凹透镜};
\end{tikzpicture}

{\small 近视眼使平行光线聚焦于视网膜前方（红
点）；凹透镜可使焦点落在视网膜上。}
\end{omfigure}

\begin{definition}[放大镜]\label{def:g11:lenses-and-eye:magnifier}
称\emph{放大镜}\index{放大镜}为焦距较短的
凸透镜，使用时物体位于\emph{焦距之内}，所成的
像为虚像、正立、放大的像（\cref{def:g11:lenses-and-eye:images}），
且人眼可在较远处舒适地观察它。
\end{definition}

\begin{example}[珠宝匠的放大镜]\label{ex:g11:lenses-and-eye:magnifier}
放大镜焦距 $f' = \qty{5.0}{cm}$，宝石位于 \qty{4.0}{cm} 处：
$1/\overline{OA'} = 1/5.0 - 1/4.0 = -1/20$，于是
$\overline{OA'} = \qty{-20}{cm}$，$\gamma = (-20)/(-4.0) = +5.0$，像为
虚像、正立，放大五倍，位于距透镜舒适的
\qty{20}{cm} 处。
\end{example}

\section{练习}

\begin{exercise}[$\star$]\label{exo:g11:lenses-and-eye:1}
计算焦距分别为 $f' = \qty{50}{cm}$、$\qty{-20}{cm}$ 和
$\qty{4.0}{mm}$ 的透镜的光焦度，并求 $+8.0$~D 透镜的焦距。
\end{exercise}

\begin{exercise}[$\star$]\label{exo:g11:lenses-and-eye:2}
一透镜的边缘比中央厚。它是会聚透镜还是
发散透镜？它能否用作放大镜？手持该透镜
置于印刷页面上方时，你看到什么？
\end{exercise}

\begin{exercise}[$\star$]\label{exo:g11:lenses-and-eye:3}
一凸透镜的焦距 $f' = \qty{10}{cm}$，物体位于其前方
\qty{20}{cm} 处。用三条特殊光线中的两条作出像，
再用共轭关系式检验像的位置与大小。
\end{exercise}

\begin{exercise}[$\star$]\label{exo:g11:lenses-and-eye:4}
一物体位于一凸透镜前方 \qty{24}{cm} 处，该凸透镜的
$f' = \qty{8.0}{cm}$。求 $\overline{OA'}$ 与 $\gamma$，并
描述像的性质、正倒与大小。
\end{exercise}

\begin{exercise}[$\star$]\label{exo:g11:lenses-and-eye:5}
在眼中：什么结构相当于凸透镜？什么起屏幕的作用？
调节时何者改变，何者不变？定义
标准眼的近点。
\end{exercise}

\begin{exercise}[$\star\star$]\label{exo:g11:lenses-and-eye:6}
一照相机镜头焦距 $f' = \qty{50}{mm}$。对准远景调焦时，
感光元件位于何处？求被摄物体位于距镜头
\qty{3.0}{m} 处时镜头与感光元件之间的距离，以及两种调焦状态之间感光元件的移动距离。
\end{exercise}

\begin{exercise}[$\star\star$]\label{exo:g11:lenses-and-eye:7}
一放大镜的焦距 $f' = \qty{6.0}{cm}$；一枚邮票位于其下方
\qty{5.0}{cm} 处。求像的位置、性质、正倒和放大率。
\end{exercise}

\begin{exercise}[$\star\star$]\label{exo:g11:lenses-and-eye:8}
一物体与一屏幕彼此相对，二者之间放置一凸透镜，
当二者距透镜均为 \qty{40}{cm} 时，屏上得到清晰的像。求
$f'$ 与 $\gamma$。为什么这一对称成像必定是等大的？
\end{exercise}

\begin{exercise}[$\star\star$]\label{exo:g11:lenses-and-eye:9}
一物体位于一凹透镜前方 \qty{50}{cm} 处，该凹透镜的
$f' = \qty{-25}{cm}$。求 $\overline{OA'}$ 与 $\gamma$，并描述
像。为什么近视者的眼睛透过其眼镜
看起来会稍小一些？
\end{exercise}

\begin{exercise}[$\star\star$]\label{exo:g11:lenses-and-eye:10}
近视眼的远点位于 \qty{40}{cm} 处。矫正透镜应具有
多大的光焦度？它把远山的像成在何处？
此时为何看远山毫不费力？
\end{exercise}

\begin{exercise}[$\star\star$]\label{exo:g11:lenses-and-eye:11}
一位老花眼读者的近点位于 \qty{80}{cm} 处，他想在
\qty{25}{cm} 处阅读：眼镜须将位于 \qty{25}{cm} 处的书页成像于
\qty{80}{cm} 处的虚页面上。求所需的光焦度。
\end{exercise}

\begin{exercise}[$\star\star\star$]\label{exo:g11:lenses-and-eye:12}
一台投影仪须对幻灯片成像，使其放大率 $\gamma = -50$，并将像投射到
距幻灯片 \qty{5.1}{m} 的屏幕上。
\begin{enumerate}
  \item 求 $\overline{OA}$、$\overline{OA'}$，已知
        $\overline{OA'} - \overline{OA} = \qty{5.1}{m}$。
  \item 由此推出投影透镜的焦距。
  \item 高 \qty{24}{mm} 的幻灯片，其像的高度是多少？
\end{enumerate}
\end{exercise}

\begin{exercise}[$\star\star\star$]\label{exo:g11:lenses-and-eye:13}
视网膜位于眼睛晶状体后方 \qty{17}{mm} 处。
\begin{enumerate}
  \item 求眼睛对无穷远处成像时的光焦度，再求对
        标准近点（\qty{25}{cm}）成像时的光焦度，由此推出
        标准眼的调节振幅。
  \item 某只年老的眼睛，近点位于 \qty{1.0}{m} 处。求其
        调节振幅，以及剩下的部分占标准调节振幅的比例。
\end{enumerate}
\end{exercise}

\begin{exercise}[$\star\star\star$]\label{exo:g11:lenses-and-eye:14}
一个标称为“$\times 3$”的放大镜，其焦距由
$3 = 0.25/f'$（$f'$ 以米计）确定。
\begin{enumerate}
  \item 求 $f'$。
  \item 物体应置于何处，才能使其虚像位于透镜前
        方 \qty{25}{cm} 处？求 $\gamma$，并将其与
        标称的 $\times 3$ 相比较。
\end{enumerate}
\end{exercise}

\begin{exercise}[$\star\star\star$]\label{exo:g11:lenses-and-eye:15}
一实物位于凸透镜前方距离 $d > 0$ 处
（$\overline{OA} = -d$，$f' > 0$）。证明：
$\overline{OA'} = \frac{f'd}{d - f'}$，并
$\gamma = \frac{f'}{f' - d}$。由此可知：当 $d > f'$ 时，像为实像、倒立；
像为虚像、正立、\emph{放大}（$\gamma > 1$），
对应 $0 < d < f'$ 的情形——只要距离足够近，每个凸透镜都可用作放大镜。
\end{exercise}

\section{题目：一个公式，三种仪器}

\begin{problem}[{周末问题——同一共轭关系
可用于聚焦照相机、配制老花镜、设计放大镜，
并测量你自己眼睛内置的变焦能力}]
\label{pb:g11:lenses-and-eye:1}
一位摄影师调焦 \qty{50}{mm} 的镜头，一位奶奶把
报纸拿在手臂伸直的距离处，一位珠宝商俯身察看一颗钻石：一行
代数式 $1/\overline{OA'} - 1/\overline{OA} = 1/f'$
（\cref{thm:g11:lenses-and-eye:conjugate}），便能胜任这三种工作——
再看人眼本身。

\medskip
\noindent\textbf{第一部分——一个公式。}一凸透镜，其
$f' = \qty{20}{cm}$。
\begin{enumerate}
  \item 物体位于 \qty{60}{cm} 处：求 $\overline{OA'}$ 与
        $\gamma$，并描述像。
  \item 将物体移至 \qty{30}{cm} 处，重新计算。与第 1 题相比，
        交换了什么？（光路是可逆的。）
  \item 将物体移至 \qty{10}{cm} 处，位于 $f'$ 以内：重新计算并描述。
  \item 总结：物体从很远处向 $F$ 靠近时，
        再越过它之后，像如何移动、性质如何变化？
  \item 证明：极远处的物体（$1/\overline{OA} \approx 0$）
        总成像于焦平面上：$\overline{OA'} = f'$。
\end{enumerate}

\medskip
\noindent\textbf{第二部分——摄影师。}
一台照相机镜头，其 $f' = \qty{50}{mm}$；感光元件高 \qty{24}{mm}，
可相对于镜头移动。
\begin{enumerate}[resume]
  \item 对远处风景调焦时，感光元件位于何处？
  \item 若一张人脸距镜头 \qty{1.0}{m}，求镜头与感光元件之间的新距离，
        以及相对于远景调焦时的移动量。
  \item 该机构允许的镜头与感光元件之间的距离至多为
        \qty{55}{mm}。求最小对焦距离。
  \item 在该最近对焦位置，求 $\gamma$。一张
        \qty{20}{cm} 的人脸，其像能否容纳于感光元件之内？
  \item 等大“微距”意味着 $\gamma = -1$，此时物体与
        感光元件应分别位于何处？为何这需要特殊镜头？
\end{enumerate}

\medskip
\noindent\textbf{第三部分——奶奶。}
随着年龄增长，近点后移：奶奶的近点位于 \qty{1.0}{m} 处。
\begin{enumerate}[resume]
  \item 解释为何报纸要拿在手臂伸直的距离处。
  \item 要在 \qty{25}{cm} 处阅读，她的眼镜须将位于
        \qty{25}{cm} 处的书页成像为位于 \qty{1.0}{m} 处的虚像，求
        所需的光焦度。
  \item 求 $\gamma$。像大了四倍，但看上去
        并不更大：请比较大小\emph{和}距离来解释。
  \item 证明：印刷字迹仅在 \qty{25}{cm} 与
        $f' \approx \qty{33}{cm}$ 之间方才清晰。那么，透过眼镜看一个\emph{远处}
        的物体又会如何？
  \item 用一句话推出为何存在双光透镜。
\end{enumerate}

\medskip
\noindent\textbf{第四部分——珠宝商，以及眼睛自身的变焦。}
珠宝商的放大镜，其 $f' = \qty{5.0}{cm}$；他的近点位于
\qty{25}{cm} 处；他的视网膜位于其眼睛晶状体后方 \qty{17}{mm} 处。
\begin{enumerate}[resume]
  \item 一块宝石位于放大镜前方 \qty{4.0}{cm} 处：其像的位置、性质，
        放大率为多少？
  \item 宝石应位于何处，才能使其像落在他的近点
        处？此时 $\gamma$ 为多少？
  \item 为何经验丰富的珠宝商反而将宝石恰好置于 $F$
        处（像位于无穷远）？
  \item 现仅考虑眼睛：求其观看远处物体时的光焦度，
        再求观看 \qty{25}{cm} 处物体时的光焦度；
        其调节量为多少屈光度？
  \item 将两个光焦度均换算为焦距，眼睛自身的焦距
        以毫米和百分比计改变了多少？
——即照相机依靠移动完成的变焦，以及奶奶借助
        $+3$~屈光度完成的变焦？
\end{enumerate}
\end{problem}
